% 第九卷 15.2.1 节：语音识别模型谱系（跨栏，约156 mm）。
% 蓝色：显式状态对齐；黄色：对齐路径求和；红色：注意力条件生成。
% 实线箭头表示“改变了前一方法的哪个机制”，虚线表示另一来源的组成部分。
\begin{tikzpicture}[
  x=1mm,y=1mm,
  figure picture,
  base/.style={figure node,minimum width=30mm,minimum height=11mm,
    inner ysep=2pt,font=\figurefont\scriptsize},
  nst/.style={base,draw=FigureInput,fill=FigureInputFill},
  nps/.style={base,draw=FigureModel,fill=FigureCream},
  nat/.style={base,draw=FigureOutput,fill=FigureOutputFill},
  fl/.style={figure flow},
  dl/.style={figure flow,draw=FigureMuted,dash pattern=on 2pt off 1.5pt},
  el/.style={font=\figurefont\scriptsize,text=FigureMuted,inner sep=1pt,align=center}
]
  \node[nst] (hmm)   at (15,40) {HMM--GMM\\帧级状态与GMM似然};
  \node[nst] (hyb)   at (15,18) {混合HMM--DNN\\网络给出状态后验};
  \node[nps] (ctc)   at (57,40) {CTC\\对单调路径求和};
  \node[nps] (rnnt)  at (57,18) {RNN-T\\路径求和＋输出历史};
  \node[nps] (chunk) at (57,-4) {分块流式编码器\\限制右侧可见范围};
  \node[nat] (aed)   at (99,40) {注意力编码器—解码器\\读取历史，无单调约束};
  \node[nat] (joint) at (99,18) {联合CTC/注意力\\CTC分支补单调约束};
  \node[nat] (wsp)   at (141,40) {Whisper\\弱监督与任务词元};
  \node[nat] (wss)   at (141,18) {Whisper-Streaming\\稳定前缀提交};

  \draw[fl] (hmm) -- node[el,left] {神经网络\\替换GMM} (hyb);
  \draw[fl] (hyb.north east) -- node[el,below right=-1mm and 0mm,pos=.45] {去掉帧级\\对齐标签} (ctc.south west);
  \draw[fl] (ctc) -- node[el,above] {引入\\输出历史} (aed);
  \draw[fl] (ctc) -- node[el,right] {加入\\预测网络} (rnnt);
  \draw[dl] (aed.south west) -- (rnnt.north east);
  \draw[fl] (aed) -- node[el,right] {共享编码器\\加CTC分支} (joint);
  \draw[fl] (aed) -- node[el,above] {扩大\\数据覆盖} (wsp);
  \draw[fl] (wsp) -- node[el,right] {缓冲重识别} (wss);
  \draw[fl] (rnnt) -- node[el,right] {因果化\\与分块} (chunk);

  % 图例
  \node[nst,minimum width=6mm,minimum height=4mm] (l1) at (88,-4) {};
  \node[anchor=west,font=\figurefont\scriptsize] at (l1.east) {显式状态对齐};
  \node[nps,minimum width=6mm,minimum height=4mm] (l2) at (111,-4) {};
  \node[anchor=west,font=\figurefont\scriptsize] at (l2.east) {对齐路径求和};
  \node[nat,minimum width=6mm,minimum height=4mm] (l3) at (134,-4) {};
  \node[anchor=west,font=\figurefont\scriptsize] at (l3.east) {注意力条件生成};
\end{tikzpicture}
